\chapter{Risks, Timeline and Publication Strategy}
\label{ch:plan}

This chapter turns the plan into a schedule. It lists the risks and how each
will be handled. It then sets the timeline and chooses target journals,
anticipates what reviewers will ask, and outlines the manuscript.

\section{Risk register}

Table~\ref{tab:risks} lists the ten risks most likely to delay or weaken the
work, with the response planned for each.

\begin{table}[htbp]
\centering
\caption[Risk register]{Risk register. Likelihood and impact are rated low,
medium or high.}
\label{tab:risks}
\small
\begin{tabularx}{\textwidth}{@{}P{3.6cm}P{1.25cm}P{1.15cm}L@{}}
\hd{Risk} & \hd{Likelihood} & \hd{Impact} & \hd{Mitigation} \\ \gap
RGBT-Tiny access is delayed & medium & high & request on day one; begin Stage B on VT-VOD50 and UVT-VOD2024 \\
\texttt{mamba\_ssm} fails on T4 & medium & medium & pure-PyTorch scan; sequences here are short enough \\
Kaggle quota is insufficient & medium & medium & two T4s per session; university cluster or Colab as contingency; fewer single-seed ablations \\
H1 is refuted & medium & medium & report it; the robustness protocol and H3 still stand as contributions \\
a competing paper appears & medium & high & monthly prior-art check; sharpen the claim toward robustness and deployment \\
motion estimation fails on thermal & medium & low & learned estimator; state reset on failure \\
scan does not export to TensorRT & high & medium & single-step state update as plain operations; plugin as last resort \\
label noise in benchmarks & medium & low & audit in Notebook~03; report on corrected labels where available \\
overfitting to one benchmark & medium & medium & cross-dataset evaluation on four further datasets \\
schedule slip & high & medium & each phase is a usable deliverable; Phase~4 can be dropped \\
\end{tabularx}
\end{table}

\section{Timeline}

The schedule assumes work begins in October 2026 and targets submission in
July 2027 (Figure~\ref{fig:gantt}). Five milestones mark the exit tests of
the phases.

\begin{figure}[htbp]
\centering
\begin{ganttchart}[
    hgrid=false, vgrid=false,
    x unit=0.98cm, y unit title=0.62cm, y unit chart=0.56cm,
    title/.style={draw=none, fill=none},
    title label font=\scriptsize\color{muted},
    title height=0.9,
    bar/.style={draw=none, fill=cTime!45, rounded corners=1.5pt},
    bar height=0.5, bar top shift=0.25,
    bar label font=\scriptsize\color{ink},
    bar label node/.append style={align=right},
    group/.style={draw=none, fill=cNeutral!55},
    group height=0.22, group top shift=0.4, group peaks height=0, group peaks width=0,
    group label font=\scriptsize\scshape\color{muted},
    canvas/.style={draw=none, fill=none},
    milestone/.style={shape=diamond, inner sep=0pt, draw=none, fill=cAlert},
    milestone label font=\scriptsize\itshape\color{cAlert},
    milestone height=0.5, milestone top shift=0.25,
    link/.style={-{Stealth[length=1.6mm]}, draw=cNeutral, line width=0.45pt},
  ]{1}{10}
  \gantttitle{2026}{3} \gantttitle{2027}{7} \\
  \gantttitle{Oct}{1} \gantttitle{Nov}{1} \gantttitle{Dec}{1} \gantttitle{Jan}{1}
  \gantttitle{Feb}{1} \gantttitle{Mar}{1} \gantttitle{Apr}{1} \gantttitle{May}{1}
  \gantttitle{Jun}{1} \gantttitle{Jul}{1} \\

  \ganttgroup{Setup}{1}{2} \\
  \ganttbar[bar/.append style={fill=cNeutral!40}]{Phase 0, environment}{1}{1} \\
  \ganttbar[bar/.append style={fill=cData!45}]{data requests and conversion}{1}{2} \\
  \ganttmilestone{M1 environment ready}{1} \\

  \ganttgroup{Still images}{2}{4} \\
  \ganttbar[bar/.append style={fill=cFuse!45}]{baselines}{2}{3} \\
  \ganttbar[bar/.append style={fill=cFuse!60}]{\method, Stage A, ablations}{3}{4} \\
  \ganttmilestone{M2 Phase 1 exit}{4} \\

  \ganttgroup{Video}{5}{7} \\
  \ganttbar{\methodS, Stage B}{5}{6} \\
  \ganttbar[bar/.append style={fill=cTime!30}]{robustness and tracking}{6}{7} \\
  \ganttmilestone{M3 H1 to H4 answered}{6} \\

  \ganttgroup{Live}{7}{9} \\
  \ganttbar[bar/.append style={fill=cRGB!45}]{export and live system}{7}{8} \\
  \ganttbar[bar/.append style={fill=cRGB!22}]{Phase 4, optional}{8}{9} \\
  \ganttmilestone{M4 live demonstration}{8} \\

  \ganttgroup{Paper}{7}{10} \\
  \ganttbar[bar/.append style={fill=cAlert!35}]{writing}{7}{9} \\
  \ganttbar[bar/.append style={fill=cAlert!25}]{prior-art check, review}{9}{9} \\
  \ganttmilestone{M5 submission}{10}
\end{ganttchart}

\vspace{2mm}
\caption[Timeline from October 2026 to July 2027]{Timeline from October 2026
to July 2027. Grey lines group the work by phase, and red diamonds mark the
five milestones. Writing begins while the live system is built, so that the
method and results sections are drafted as soon as the hypotheses are
answered.}
\label{fig:gantt}
\end{figure}


\section{Target journals}
\label{sec:journals}

Table~\ref{tab:journals} lists suitable journals with their latest
reported impact factors, which must be confirmed in Journal Citation Reports
before submission. The work can support up to three papers:
\begin{itemize}
  \item \emph{Paper A}, the method and the robustness protocol. It is the
        main paper and suits IEEE TCSVT, Pattern Recognition or Information
        Fusion.
  \item \emph{Paper B}, the deployed real-time system with edge
        measurements. It suits the Journal of Real-Time Image Processing
        (Springer) or Engineering Applications of Artificial Intelligence.
  \item \emph{Paper C}, an extended version for remote sensing, if the
        drone results are strong. It suits IEEE TGRS or the International
        Journal of Remote Sensing (Taylor \& Francis).
\end{itemize}

\begin{table}[htbp]
\centering
\caption[Candidate journals]{Candidate journals. Impact factors are as
reported by secondary sources for the 2025 release and are marked
unverified.}
\label{tab:journals}
\small
\begin{tabularx}{\textwidth}{@{}LP{2.1cm}R{1.35cm}P{4.2cm}@{}}
\hd{Journal} & \hd{Publisher} & \hd{JIF\unv} & \hd{Fit} \\ \gap
Information Fusion & Elsevier & 15.5 & Paper A if framed around fusion robustness \\
IEEE Trans.\ Image Processing & IEEE & 13.7 & Paper A with strong analysis; demanding \\
IEEE Trans.\ Circuits Syst.\ Video Technol. & IEEE & 11.1 & \textbf{Paper A, best fit}: video and streaming \\
IEEE Trans.\ Multimedia & IEEE & 9.7 & Paper A; published Fusion-Mamba \\
IEEE Trans.\ Geosci.\ Remote Sensing & IEEE & 8.6 & Paper C \\
Eng.\ Appl.\ of Artificial Intelligence & Elsevier & 8.0 & Paper B \\
Pattern Recognition & Elsevier & 7.6 & Paper A \\
Neural Networks & Elsevier & 6.3 & state-space analysis \\
Neural Computing and Applications & Springer & 4.5--6 & fallback for A \\
Computer Vision and Image Understanding & Elsevier & 3.6 & published UAVDet \\
J.\ Real-Time Image Processing & Springer & 3.0 & \textbf{Paper B, best fit} \\
The Visual Computer & Springer & 2.9 & fallback \\
Int.\ J.\ Remote Sensing & Taylor \& Francis & 2.6 & Paper C \\
Machine Vision and Applications & Springer & 2.3 & fallback \\
\end{tabularx}
\end{table}

\section{What reviewers will ask}
\label{sec:reviewers}

Reviewers at these journals ask for the same things repeatedly. The protocol
of Chapter~\ref{ch:protocol} answers each in advance.
\begin{itemize}
  \item At least three datasets, including a video dataset and a dataset
        with real misalignment.
  \item Comparisons against 2025--2026 methods, not only against YOLOv8:
        WaveMamba, COMO, MambaST and the alignment-free video detectors.
  \item Parameters, GFLOPs and latency for every method, measured on the
        same GPU at batch one, end to end. Jetson numbers are needed if
        real time is claimed.
  \item Mean and standard deviation over at least three seeds, with
        confidence intervals that respect the correlation of frames within
        sequences.
  \item An ablation for every component, failure cases and visualisations
        of what the state holds.
  \item Code, weights and a data-availability statement.
\end{itemize}

\section{Manuscript outline}

Table~\ref{tab:outline} gives the planned structure of Paper~A. The method
and experiments sections take most of the space, as reviewers expect.

\begin{table}[htbp]
\centering
\caption[Outline of the main manuscript]{Outline of the main manuscript
(Paper A), with approximate lengths for a journal of the IEEE Transactions
format.}
\label{tab:outline}
\small
\begin{tabularx}{\textwidth}{@{}P{3.7cm}LR{1.4cm}@{}}
\hd{Section} & \hd{Content} & \hd{Pages} \\ \gap
Introduction & problem, gap, contributions, figure of the streaming setting & 1.25 \\
Related work & YOLO and SSM detectors; RGB--T fusion; temporal and streaming detection & 1.25 \\
Method & CMFM with gated step; TSM; motion compensation; training & 3.0 \\
Robustness protocol & perturbations, severities, event metrics & 0.75 \\
Experiments & datasets, baselines, main results, ablations, robustness, efficiency & 4.0 \\
Analysis & what the state holds, failure cases, limitations & 1.0 \\
Conclusion and statements & conclusion, ethics, data and code availability & 0.5 \\
\end{tabularx}
\end{table}

\section{Before submission}
\label{sec:submission}

Five checks come immediately before submission.
\begin{enumerate}
  \item Prior art is searched again on arXiv and in the target journals
        across the four lanes of Figure~\ref{fig:timeline}. The claims are
        adjusted to whatever has appeared.
  \item The code is released under the AGPL-3.0 licence that Ultralytics
        requires, with configurations and the run logs behind every table.
  \item Weights and the robustness protocol's corruption seeds are released,
        so that others can evaluate on identical streams.
  \item The ethics and data statements of Chapter~\ref{ch:ethics} are
        completed.
  \item A colleague who has not seen the project reproduces one main table
        from the released material.
\end{enumerate}
